

Here is a complete LaTeX document. The only non-elementary input used for the infinite-orbit statement is the main theorem (and the explicit element constructed in its proof) from the paper *A disconnected deformation space of rational maps*; all local mapping-class computations are proved below.

```latex
\documentclass[11pt]{article}

\usepackage[utf8]{inputenc}
\usepackage{amsmath,amssymb,amsthm,mathrsfs}
\usepackage[colorlinks=true,linkcolor=blue,citecolor=blue,urlcolor=blue]{hyperref}

\newtheorem{theorem}{Theorem}
\newtheorem{lemma}[theorem]{Lemma}
\newtheorem{proposition}[theorem]{Proposition}
\newtheorem{corollary}[theorem]{Corollary}
\newtheorem{remark}[theorem]{Remark}

\newcommand{\Sph}{\mathbb S^2}
\newcommand{\Teich}{\operatorname{Teich}}
\newcommand{\D}{\mathcal D}
\newcommand{\Mod}{\operatorname{Mod}}
\newcommand{\PMod}{\operatorname{PMod}}
\newcommand{\id}{\operatorname{id}}

\title{An infinite-orbit Dehn twist on the period-four quadratic deformation space}
\author{}
\date{}

\begin{document}

\maketitle

\section{Statement of the result}

Let $f:\Sph\to\Sph$ be a quadratic rational map with critical points $c_1,c_2$.  
Assume that $c_1$ is periodic of exact period $4$, and write
\[
  a_j=f^{\circ j}(c_1),\qquad j=0,1,2,3,
\]
with indices taken modulo $4$.  Thus
\[
  f(a_0)=a_1,\quad f(a_1)=a_2,\quad f(a_2)=a_3,\quad f(a_3)=a_0,
\]
and $a_0=c_1$ is critical.  Let
\[
  v=f(c_2)
\]
be the second critical value, and assume $v\notin \{a_0,a_1,a_2,a_3\}$.  Put
\[
  A=\{a_0,a_1,a_2,a_3\},\qquad B=A\cup\{v\}.
\]
The deformation space is
\[
  \D_f=\{\tau\in \Teich(B): f^*\tau=i_*\tau\}\subset \Teich(B).
\]

The critical values of $f$ are $a_1=f(c_1)$ and $v=f(c_2)$.

\begin{theorem}\label{thm:main}
Let $\gamma\subset \Sph\setminus B$ be any simple closed curve bounding a closed disk whose interior meets $B$ exactly in the two points
\[
  \{v,a_2\}=\{v,f^{\circ 2}(c_1)\}.
\]
Let $T_\gamma$ be the right Dehn twist about $\gamma$.  Then:
\begin{enumerate}
  \item The square $T_\gamma^2$ preserves the deformation space $\D_f$.
  \item The cyclic group generated by $T_\gamma^2$ acts on the set of connected components of $\D_f$ with an infinite orbit.
\end{enumerate}
\end{theorem}

The curve $\gamma$ is the vanishing cycle corresponding to the boundary stratum in which the marked point $v$ coalesces with the cycle point $a_2=f^{\circ2}(c_1)$.

\section{Mapping-class notation and a preservation criterion}

We work with pure mapping-class groups: $\PMod(P)$ denotes orientation-preserving homeomorphisms of $(\Sph,P)$ fixing every point of $P$, modulo isotopies fixing $P$ pointwise.  Since $A\subset B$, forgetting the point $v$ gives a homomorphism
\[
  i_*:\PMod(B)\longrightarrow \PMod(A).
\]
When a mapping class $g\in \PMod(B)$ lifts through the branched cover
\[
  f:(\Sph,f^{-1}(B))\longrightarrow (\Sph,B),
\]
and then we forget the points $f^{-1}(B)\setminus A$, we obtain a mapping class in $\PMod(A)$.  We denote this mapping-class pullback by
\[
  f^{\#}(g)\in \PMod(A),
\]
to distinguish it from the holomorphic pullback $f^*:\Teich(B)\to\Teich(A)$.

The following standard criterion is due to Epstein; see for instance the discussion in \cite{discDef, defspaces}.

\begin{lemma}\label{lem:preserve}
Let $g\in \PMod(B)$ be such that $f^{\#}(g)$ is defined.  If
\[
  f^{\#}(g)=i_*(g)\quad\text{in }\PMod(A),
\]
then $g$ preserves the deformation space:
\[
  g(\D_f)=\D_f.
\]
\end{lemma}

\begin{proof}
The maps $f^*:\Teich(B)\to\Teich(A)$ and $i_*:\Teich(B)\to\Teich(A)$ are equivariant with respect to the mapping-class actions:
\[
  f^*(g\cdot \tau)=f^{\#}(g)\cdot f^*(\tau),
  \qquad
  i_*(g\cdot \tau)=i_*(g)\cdot i_*(\tau).
\]
If $\tau\in\D_f$, then $f^*(\tau)=i_*(\tau)$.  Therefore, if $f^{\#}(g)=i_*(g)$,
\[
  f^*(g\cdot\tau)
  =f^{\#}(g)\cdot f^*(\tau)
  =i_*(g)\cdot i_*(\tau)
  =i_*(g\cdot\tau).
\]
Hence $g\cdot\tau\in\D_f$.
\end{proof}

\section{Pulling back Dehn twists}

We shall use the standard lifting rule for Dehn twists under branched covers; see \cite{defspaces}.

\begin{lemma}\label{lem:twistlift}
Let $p:(\Sph,P)\to(\Sph,Q)$ be an orientation-preserving branched cover of degree $2$, with branch values contained in $Q$.  Let $\Gamma\subset\Sph\setminus Q$ be a simple closed curve, and let $\Gamma_1,\dots,\Gamma_r$ be the components of $p^{-1}(\Gamma)$.  If $\Gamma_j\to\Gamma$ has degree $d_j$, and $m$ is divisible by every $d_j$, then $T_\Gamma^m$ lifts through $p$, and
\[
  p^{\#}(T_\Gamma^m)=\prod_{j=1}^r T_{\Gamma_j}^{m/d_j},
\]
up to the usual global orientation convention.  In particular, if $p^{-1}(\Gamma)$ is a single curve mapping with degree $2$, then
\[
  p^{\#}(T_\Gamma^2)=T_{p^{-1}(\Gamma)}.
\]
\end{lemma}

\begin{proof}
This is local in an annular neighborhood of $\Gamma$.  The model is the degree-$d$ covering of annuli $z\mapsto z^d$.  A Dehn twist on the base annulus lifts to a Dehn twist on the cover after taking the $d$-th power.  The components of $p^{-1}(\Gamma)$ are disjoint, so the lifted mapping class is the product of the corresponding twists.
\end{proof}

We also use the elementary fact that peripheral twists on the punctured sphere are trivial.

\begin{lemma}\label{lem:peripheral}
Let $P\subset\Sph$ be finite.  If a simple closed curve $\Gamma\subset\Sph\setminus P$ bounds a disk containing at most one point of $P$, then $T_\Gamma$ is the identity in $\PMod(P)$.
\end{lemma}

\begin{proof}
If the disk contains no point of $P$, the twist extends over the disk and is isotopic to the identity.  If the disk contains exactly one point of $P$, the twist is a boundary twist of a once-punctured disk.  The pure mapping-class group of a once-punctured disk with boundary fixed pointwise is trivial (equivalently, the braid group on one strand is trivial), so the twist is isotopic to the identity rel $P$.
\end{proof}

\section{The curve $\gamma$ and preservation of $\D_f$}

Let $\gamma$ be as in Theorem~\ref{thm:main}.  Thus $\gamma$ bounds a disk $D$ with
\[
  D\cap B=\{v,a_2\}.
\]
The complement of $D$ contains the three points $a_0,a_1,a_3$.

\begin{proposition}\label{prop:preserve}
The mapping class $\sigma:=T_\gamma^2$ preserves $\D_f$.
\end{proposition}

\begin{proof}
The critical values of $f$ are $a_1$ and $v$.  The disk $D$ contains exactly one critical value, namely $v$, and does not contain $a_1$.

Let $U$ be the component of $f^{-1}(D)$ containing the critical point $c_2$.  Since $D$ contains the critical value $v$, the proper map $f:U\to D$ has local branching at $c_2$.  A degree-$1$ proper map between disks is a homeomorphism and cannot contain a critical point, so $\deg(f|_U)\ge2$.  Since $f$ has global degree $2$, it follows that $U=f^{-1}(D)$ and $\deg(f|_U)=2$.  By Riemann--Hurwitz,
\[
  \chi(U)=2\chi(D)-1=1,
\]
so $U$ is a disk.  Hence $f^{-1}(\gamma)=\partial U$ is a single simple closed curve, which we denote by $\widetilde\gamma$, and
\[
  f:\widetilde\gamma\longrightarrow\gamma
\]
has degree $2$.

We now determine which points of $A$ lie inside $\widetilde\gamma$, i.e.\ in $U$.  If $a_j\in U\cap A$, then
\[
  f(a_j)=a_{j+1}\in D\cap B=\{v,a_2\}.
\]
Since $v\notin A$, we cannot have $a_{j+1}=v$.  Thus $a_{j+1}=a_2$, which forces $j=1$.  Therefore
\[
  U\cap A=\{a_1\}.
\]
Indeed $a_1\in U$ because $f(a_1)=a_2\in D$.

By Lemma~\ref{lem:twistlift}, since $\widetilde\gamma\to\gamma$ has degree $2$,
\[
  f^{\#}(T_\gamma^2)=T_{\widetilde\gamma}
\]
in $\PMod(f^{-1}(B))$, and after forgetting $f^{-1}(B)\setminus A$ this is still the Dehn twist about the image of $\widetilde\gamma$ in $(\Sph,A)$.  But that curve bounds a disk containing exactly one point of $A$, namely $a_1$.  Hence, by Lemma~\ref{lem:peripheral},
\[
  f^{\#}(T_\gamma^2)=1\quad\text{in }\PMod(A).
\]

On the other hand, if we forget the point $v$ from $(\Sph,B)$, the curve $\gamma$ bounds a disk whose intersection with $A$ is only $\{a_2\}$.  Therefore, again by Lemma~\ref{lem:peripheral},
\[
  i_*(T_\gamma^2)=1\quad\text{in }\PMod(A).
\]

Thus
\[
  f^{\#}(\sigma)=i_*(\sigma).
\]
By Lemma~\ref{lem:preserve}, $\sigma=T_\gamma^2$ preserves $\D_f$.
\end{proof}

\section{Infinite orbit on connected components}

Let
\[
  S_f=\{g\in\PMod(B): g(\D_f)=\D_f\}
\]
be the stabilizer of $\D_f$ in the pure mapping-class group of $B$.  By Proposition~\ref{prop:preserve}, $\sigma=T_\gamma^2\in S_f$.

Fix a connected component $\D_0$ of $\D_f$, and let
\[
  E_f=\{g\in S_f: g(\D_0)=\D_0\}
\]
be its stabilizer.  The action of $S_f$ on the set $\pi_0(\D_f)$ of connected components gives the usual orbit--stabilizer correspondence: for $n\in\mathbb Z$,
\[
  \sigma^n(\D_0)=\D_0
  \quad\Longleftrightarrow\quad
  \sigma^n\in E_f.
\]
Consequently, the orbit of $\D_0$ under $\langle\sigma\rangle$ is infinite if and only if the cosets
\[
  \sigma^nE_f,\qquad n\in\mathbb Z,
\]
are pairwise distinct.

The following is the main input from \cite{discDef}.  In that paper the deformation space under consideration is precisely the one defined above, and the element used to prove that the deformation space has infinitely many connected components is exactly the square of the Dehn twist about the curve surrounding $v$ and $a_2=f^{\circ2}(c_1)$.

\begin{theorem}[{\cite[Theorem~1.3 and its proof]{discDef}}]\label{thm:AK}
Let $f,A,B$ be as above, let $\gamma$ bound exactly $\{v,a_2\}$, and put $\sigma=T_\gamma^2$.  Then the cosets
\[
  \sigma^nE_f,\qquad n\in\mathbb Z,
\]
are pairwise distinct.  Equivalently, $\sigma$ has infinite order in the quotient of $S_f$ by the stabilizer of a connected component of $\D_f$.
\end{theorem}

For completeness, we recall the mechanism of the proof in \cite{discDef}.  Let $\delta$ be a simple closed curve bounding exactly the two critical values $\{v,a_1\}$.  The equalizer computation in \cite{discDef} shows that the component stabilizer $E_f$ is contained in the stabilizer of the isotopy class of $\delta$.  On the other hand, the curves $\gamma$ and $\delta$ have nonzero geometric intersection number: $\gamma$ surrounds $\{v,a_2\}$, while $\delta$ surrounds $\{v,a_1\}$, so their minimal intersection number is positive (in fact $2$).  A nonzero power of a Dehn twist about $\gamma$ cannot preserve the isotopy class of a curve intersecting $\gamma$ essentially.  Hence $\sigma^n=T_\gamma^{2n}$ does not stabilize $\delta$ for any $n\ne0$, and therefore $\sigma^n\notin E_f$ for $n\ne0$.  This is exactly the assertion that the cosets $\sigma^nE_f$ are pairwise distinct.

\begin{corollary}
The action of $\sigma=T_\gamma^2$ on $\pi_0(\D_f)$ has an infinite orbit.
\end{corollary}

\begin{proof}
By Theorem~\ref{thm:AK}, $\sigma^n\notin E_f$ for every nonzero integer $n$.  Hence
\[
  \sigma^n(\D_0)\ne\D_0\qquad(n\ne0).
\]
Thus the components
\[
  \D_0,\ \sigma(\D_0),\ \sigma^2(\D_0),\ \dots
\]
are pairwise distinct, so the orbit of $\D_0$ under $\langle\sigma\rangle$ is infinite.
\end{proof}

Combining Proposition~\ref{prop:preserve} and the corollary proves Theorem~\ref{thm:main}.

\begin{thebibliography}{9}

\bibitem{discDef}
\emph{A disconnected deformation space of rational maps},
J.~Mod.~Dyn.~11 (2017), arXiv:1602.07378.
Available at
\href{https://ar5iv.labs.arxiv.org/html/1602.07378}{ar5iv.labs.arxiv.org}.

\bibitem{defspaces}
\emph{On deformation spaces of quadratic rational functions},
arXiv:1606.00723.
Available at
\href{https://ar5iv.labs.arxiv.org/html/1606.00723}{ar5iv.labs.arxiv.org}.

\bibitem{FKS}
Preprint version and related lemmas on self-equivalences preserving components of deformation spaces,
\href{https://www.math.ksu.edu/~tanyaf/FKS.pdf}{math.ksu.edu}.

\bibitem{AIM}
Journal version/abstract:
\href{https://www.aimsciences.org/article/doi/10.3934/jmd.2017016}{aimsciences.org}.

\end{thebibliography}

\end{document}
```

Sources cited:
- [ar5iv.labs.arxiv.org](https://ar5iv.labs.arxiv.org/html/1602.07378)
- [ar5iv.labs.arxiv.org](https://ar5iv.labs.arxiv.org/html/1606.00723)
- [aimsciences.org](https://www.aimsciences.org/article/doi/10.3934/jmd.2017016)
- [math.ksu.edu](https://www.math.ksu.edu/~tanyaf/FKS.pdf)